% =========================================================
\section{Conclusion and discussion}
\label{sec:conclusion}
% =========================================================

We studied the first interaction between a smooth flip bifurcation
and the endogenous switching geometry generated by radial gradient
clipping.  The mechanism operates on two sharply separated scales:
clipping is locally invisible near the critical point and becomes
dynamically relevant only at finite amplitude.

\paragraph{Local scale.}
For every fixed $c>0$ the clipped map $F_{\eta,c}$ coincides with the
ordinary gradient map $G_\eta$ on a neighborhood of a nondegenerate
critical point; hence the fixed point, its linearization, its local
invariant structures, its center-manifold reduction, and its smooth
flip coefficient are independent of the clipping threshold.  In
particular, for a strict local minimum $x_*$ with simple maximal
Hessian eigenvalue $\lambda_*$, the first stability threshold remains
\begin{equation}
\eta_f
=
\frac{2}{\lambda_*}.
\label{eq:conclusion-eta-f}
\end{equation}
Under the supercriticality condition
\begin{equation}
\Gamma_*
=
3T_*[v_*,v_*,H_*^{-1}g_*]
-
Q_*[v_*^4]
>0,
\label{eq:conclusion-Gamma}
\end{equation}
the loss of stability creates a stable smooth period-two branch
\begin{equation}
x_\pm(\mu)
=
x_*
\pm
A_*\sqrt{\mu}\,v_*
+
O(\mu),
\qquad
\mu=\eta-\eta_f,
\qquad
A_*^2
=
\frac{3\lambda_*^2}{\Gamma_*}.
\label{eq:conclusion-amplitude}
\end{equation}

\paragraph{Finite-amplitude scale.}
Clipping becomes active only when this orbit reaches the switching
hypersurface
\[
\Sigma_c
=
\{x:\|\nabla V(x)\|_\diamond=c\}.
\]
The gradient structure makes the interaction rigid: every smooth
period-two orbit satisfies the exact identity
\begin{equation}
\nabla V(x_+)
=
-\nabla V(x_-),
\label{eq:conclusion-gradient-balance}
\end{equation}
so both points have identical clipping exposure and reach $\Sigma_c$
simultaneously.  Matching the square-root amplitude scale with the
threshold produces the principal asymptotic law
\begin{equation}
\boxed{
\eta_{\mathrm{sc}}(c)-\eta_f
=
\frac{\Gamma_*}{
3\lambda_*^4
\|v_*\|_\diamond^2
}
c^2
+
O(c^4),
}
\label{eq:conclusion-main-law}
\end{equation}
which for Euclidean clipping becomes
\begin{equation}
\boxed{
\eta_{\mathrm{sc}}(c)
=
\frac{2}{\lambda_*}
+
\frac{\Gamma_*}{
3\lambda_*^4
}
c^2
+
O(c^4).
}
\label{eq:conclusion-euclidean}
\end{equation}
The even-power expansion reflects an additional symmetry of the
gradient problem: in the signed square-root parameter
$\varepsilon=\sqrt{\eta-\eta_f}$ the half-difference of the two-cycle
is odd while its midpoint is even, which together with
\eqref{eq:conclusion-gradient-balance} removes the cubic correction
from the contact curve.

\paragraph{Post-contact scale.}
A clipped period-two orbit cannot mix smooth and clipped points, so
any continuation past the first contact enters the clipped regime at
both points simultaneously, and every fully clipped period-two orbit
obeys the exact step-length constraint
\begin{equation}
\|y_+-y_-\|_\diamond
=
\eta c.
\label{eq:conclusion-step-law}
\end{equation}
Most importantly, radial clipping creates an unavoidable neutral
direction: the monodromy matrix
\[
M_c
=
DF_{\eta,c}(y_-)
DF_{\eta,c}(y_+)
\]
of a fully clipped period-two orbit satisfies
\begin{equation}
\boxed{
1\in\sigma(M_c).
}
\label{eq:conclusion-unit}
\end{equation}
Every fully clipped two-cycle is therefore nonhyperbolic: the radial
normalization removes first-order information in the
gradient-magnitude direction.  The solvable planar model of
Section~\ref{sec:planar-model} shows that this unit multiplier is not
a formal artifact --- it corresponds to an actual loss of isolation of
the periodic orbit.

% ---------------------------------------------------------
\subsection{Dynamical interpretation}
\label{subsec:conclusion-interpretation}
% ---------------------------------------------------------

The results separate two effects that are often conflated.  The
first is the loss of local stability of the minimum: a smooth
gradient-descent phenomenon governed by $D^2V(x_*)$, $D^3V(x_*)$,
$D^4V(x_*)$, and unaffected by clipping.  The second is the
modification of an already finite-amplitude trajectory once it
reaches $\|\nabla V\|_\diamond=c$.  Clipping does not postpone the
first local flip threshold; it introduces a second dynamical scale
beyond it, the switching-contact scale
$\eta_{\mathrm{sc}}(c)-\eta_f=O(c^2)$.  This distinction is useful
when interpreting large-step optimization dynamics: linear stability
determines when the critical point loses stability, while the
clipping threshold determines when the resulting nonlinear oscillation
begins to experience a different piecewise-smooth evolution.

% ---------------------------------------------------------
\subsection{Limitations}
\label{subsec:conclusion-limitations}
% ---------------------------------------------------------

The theory focuses on an isolated nondegenerate minimum with a simple
maximal Hessian eigenvalue, which isolates a one-dimensional flip mode
but excludes the additional neutral directions that arise when the set
of minimizers forms a manifold.  Second, the post-contact unit
multiplier prevents a universal hyperbolic continuation theorem:
under the present assumptions alone, the existence, uniqueness, and
nonlinear stability of a fully clipped continuation are not
determined.  Third, the paper treats radial clipping; coordinatewise
clipping, adaptive normalization, momentum, and stochastic mini-batch
updates generate different switching geometries.  Finally, the
analysis is local in the bifurcation parameter and does not classify
the remote global dynamics after the orbit enters the clipped region.

% ---------------------------------------------------------
\subsection{Directions for further work}
\label{subsec:conclusion-future}
% ---------------------------------------------------------

Several extensions appear natural: reducing the nonhyperbolic
post-contact dynamics to a higher-order scalar or center-manifold
normal form, which the exact unit multiplier suggests may form a new
class of degenerate border interactions specific to radial gradient
maps; replacing the isolated minimum by a Morse--Bott manifold of
minima, combining the $-1$ flip direction with $+1$ tangent
directions; degenerate smooth flips with $\Gamma_*=0$, where a
matching principle $r(\mu)\asymp\mu^\gamma$, $d(c)\asymp c^p$
suggests the general contact scale
$\mu_{\mathrm{sc}}\asymp c^{p/\gamma}$; and determining which
structures persist for momentum, stochastic gradients, and nonradial
clipping.

% ---------------------------------------------------------
\subsection{Final perspective}
\label{subsec:conclusion-final-perspective}
% ---------------------------------------------------------

Near a nondegenerate minimum, radial gradient clipping is locally
invisible: $F_{\eta,c}=G_\eta$, so the smooth map determines when the
first oscillatory instability is born.  Once the oscillation reaches
finite amplitude, the same potential generates an endogenous switching
boundary at the explicit scale
\[
\eta_{\mathrm{sc}}(c)-\eta_f
\sim
\frac{\Gamma_*}{
3\lambda_*^4
\|v_*\|_\diamond^2
}
c^2.
\]
At that point the dynamics cease to be governed solely by the smooth
flip normal form: radial clipping removes the gradient-magnitude
degree of freedom and forces the post-contact period-two dynamics to
carry an exact neutral multiplier.  The clipped gradient map thus
exhibits a sharp transition from smooth bifurcation theory to
intrinsically nonhyperbolic switching dynamics, with the transition
scale determined explicitly by the local geometry of the underlying
potential.
